\documentclass[11pt,a4paper]{article}
\usepackage[utf8]{inputenc}
\usepackage[T1]{fontenc}
\usepackage[romanian]{babel}
\usepackage[margin=2cm]{geometry}
\usepackage{array}
\usepackage{tabularx}
\usepackage{enumitem}
\usepackage{graphicx}
\usepackage{setspace}
\usepackage{xcolor}
\usepackage{hyperref}
\usepackage{amssymb}
\setlength{\parindent}{0pt}
\newcommand{\fillline}[1][6cm]{\underline{\hspace{#1}}}
\newcommand{\fillbox}[1][0.4cm]{\framebox[#1]{\rule{0pt}{#1}}}
\newcommand{\checkboxempty}{$\square$}

\begin{document}

\begin{tabular}{p{5cm} p{5cm} p{5cm}}
\textbf{PRIMĂRIA} & \textbf{*802014*} & Nr. \fillline[3cm] / \fillline[1.5cm] \\
\textbf{CLUJ-NAPOCA} & & \\
\end{tabular}

\vspace{0.5cm}

Către,

\begin{center}
\textbf{\underline{DIRECŢIA DE ASISTENŢĂ SOCIALĂ ŞI MEDICALĂ}}\\
\underline{Serviciul Asistenţa Persoanelor cu Dizabilităţi}
\end{center}

\vspace{0.5cm}

\hspace{1cm}Subsemnatul(a) \underline{ {{nume_titular}} } -- CNP \underline{ {{cnp_titular}} } domiciliat(ă) în Cluj-Napoca, str. \fillline[8cm], nr. \fillline[1.5cm], ap. \fillline[1.5cm], sunt persoana cu handicap/reprezentant legal al persoanei cu handicap \fillline[6cm] cu domiciliul în Cluj-Napoca, str. \fillline[5cm], nr. \fillline[1.5cm], ap. \fillline[1.5cm], solicit luarea în evidenţă în vederea eliberării \textbf{abonamentului pentru transportul urban}.

\vspace{0.2cm}

\hspace{1cm}Grad de handicap: \underline{ {{grad_handicap}} }.

\vspace{0.2cm}

\hspace{1cm}Tip abonament solicitat: \underline{ {{tip_abonament}} }.

\vspace{0.4cm}

De asemenea solicit:

\vspace{0.2cm}

Un abonament nominal asistent personal

\hspace{0.5cm}Nume, prenume \fillline[6cm] adresa \fillline[7cm]

\vspace{0.2cm}

Un abonament nenominal însoţitor.

\vspace{0.3cm}

\hspace{0.5cm}Anexez prezentei:$^{1}$ -- certificat de incadrare in grad de handicap;\\
\hspace*{3cm} -- act de identitate persoană cu handicap;\\
\hspace*{3cm} -- act de identitate reprezentant legal.

\vspace{0.4cm}

Tel. de contact \fillline[5cm]

\vspace{0.4cm}

\begin{itemize}[leftmargin=*,label=\checkboxempty]
\item Prin prezenta solicit comunicarea răspunsului pe următoarea adresă de e-mail:\\
\fillline[14cm]
\item Mă oblig să comunic instituției orice modificare intervine în legătură cu această adresă de e-mail
\item Îmi exprim consimțământul ca Primăria Municipiului Cluj-Napoca să comunice orice informații, date personale, clarificări și completări pe adresa de e-mail indicată mai sus
\item Am luat la cunoștință faptul că în cazul nefuncționării serverului de e-mail comunicat sau în cazul adresei greșite de e-mail, Municipiul Cluj-Napoca nu poate fi tras la răspundere pentru acest lucru
\end{itemize}

\vspace{1cm}

Data: \fillline[3cm] \hfill Semnătura: \fillline[4cm]

\vspace{0.8cm}

\hrule
\vspace{0.2cm}
\small
http://dasmclujnapoca.ro\\
telefon/e-mail: 0264599316 / dizabilitati@dasmclujnapoca.ro\\
Cluj-Napoca, str. Venus FN între 20-22

\vspace{0.3cm}

\footnotesize
$^{1}$Primăria municipiului Cluj-Napoca asigură gratuit fotocopierea documentelor. Pentru documentele emise de către alte instituții publice, organe de specialitate ale administrației publice centrale și locale, precum și persoane juridice de drept privat, care potrivit legii au obținut statut de utilitate publică sau sunt autorizate să presteze un serviciu public, în regim de putere publică, pe care solicitantul nu le depune, acesta are posibilitatea de a-și da consimțământul expres ca Primăria municipiului Cluj-Napoca să obțină în numele său copii sau extrase ale acestora.

\vspace{0.3cm}
\hrule
\vspace{0.2cm}
\small
www.primariaclujnapoca.ro\\
telefon: 0264/596.030\\
Cluj-Napoca, str. Moţilor, nr. 3 \hfill \textbf{*802014*}

\vspace{0.2cm}
\hrule
\vspace{0.2cm}
\footnotesize
Timp estimativ de completare: \textbf{5 minute}\\
Datele personale care vă sunt solicitate prin prezenta cerere vor fi prelucrate numai în vederea procesării și soluționării solicitării dumneavoastră. Primăria municipiului Cluj-Napoca garantează securitatea procesării datelor și arhivarea acestora în conformitate cu prevederile legale în vigoare. Responsabilul Primăriei municipiului Cluj-Napoca cu protecția datelor poate fi contactat pe adresa de email dpo@primariaclujnapoca.ro.\\
În conformitate cu Regulamentul nr. 679/2016 aveți dreptul de a solicita Primăriei municipiului Cluj-Napoca, în ceea ce priveşte datele cu caracter personal referitoare la persoana vizată, accesul la acestea, rectificarea sau ştergerea acestora sau restricţionarea prelucrării sau a dreptului de a vă opune prelucrării, precum şi a dreptului la portabilitatea datelor.

\end{document}
